\RequirePackage{luatex85}
\documentclass[border=10pt,12pt]{standalone}

\usepackage{cfr-lm}
\usepackage{amssymb}
\usepackage{pgf}
\usepackage{tikz}
\usetikzlibrary{arrows.meta,calc,positioning}

\begin{document}

\sffamily\bfseries

%% Actor x-positions (cm)
%% P=0  S=9  M=19
\def\px{0}
\def\sx{9}
\def\mx{19}
\def\ybot{34.4}

\begin{tikzpicture}[
  >=Stealth,
  yscale=-1,
]

\input{common.tex}

\tikzstyle{actor}=[rectangle, minimum width=3.6cm, minimum height=0.7cm,
  inner sep=4pt, rounded corners=2pt, align=center, font=\bfseries]
\tikzstyle{ll}=[densely dashed, gray!60, line width=0.5pt]
\tikzstyle{lldead}=[densely dashed, red!40, line width=0.5pt]
\tikzstyle{msg}=[->, line width=0.65pt]
\tikzstyle{rsp}=[->, densely dashed, line width=0.65pt]
\tikzstyle{broken}=[red!70, densely dotted, line width=0.8pt]
\tikzstyle{lbl}=[font=\small, inner sep=2pt]
\tikzstyle{note}=[rectangle, rounded corners=2pt, inner sep=5pt,
  font=\small, align=center]
\tikzstyle{seclbl}=[font=\small\itshape, text=gray!70!black, anchor=west]

%% ── actor headers ───────────────────────────────────────────────────────
\node[actor, fill=pbox, text=ptxt] (Ph) at (\px,0) {Primary};
\node[actor, fill=sbox, text=stxt] (Sh) at (\sx,0) {Secondary};
\node[actor, fill=mbox, text=mtxt] (Mh) at (\mx,0) {Monitor};

%% lifelines
\draw[ll] (\px,0.35) -- (\px,\ybot);
\draw[ll] (\sx,0.35) -- (\sx,\ybot);
\draw[ll] (\mx,0.35) -- (\mx,\ybot);

%% ── steady state: two independent signals ──────────────────────────────
\node[seclbl] at (-3.4,1.3) {steady state: two independent signals per node};

\draw[msg] (\mx,1.6) -- (\px,1.6);
\node[lbl, above] at (\mx-2.6,1.6) {health check (libpq, like \texttt{pg\_isready})};
\draw[rsp] (\px,2.2) -- (\mx,2.2)
  node[lbl,midway,above] {ok};

\draw[msg] (\px,3.4) -- (\mx,3.4);
\node[lbl, above] at (\px+4.6,3.4) {node\_active(reported=primary, pgIsRunning=true)};
\draw[rsp] (\mx,4.0) -- (\px,4.0);
\node[lbl, above] at (\mx-2.2,4.0) {goal=primary};

\draw[msg] (\sx,5.2) -- (\mx,5.2);
\node[lbl, above] at (\sx+3.8,5.2) {node\_active(reported=secondary, pgIsRunning=true)};
\draw[rsp] (\mx,5.8) -- (\sx,5.8);
\node[lbl, above] at (\mx-2.2,5.8) {goal=secondary};

\node[note, fill=gray!12, draw=gray!55, text width=15cm, font=\small]
  at (9.5, 7.1)
  {a node is only ever unhealthy from the combination of what its own
   keeper reports \emph{and} what the monitor's direct health check finds
   -- never from either alone};

%% ── primary's Postgres fails locally ────────────────────────────────────
\node[seclbl] at (-3.4,8.7)
  {primary's Postgres fails locally (keeper stays up, e.g. disk full)};

\draw[broken] (\mx,9.6) -- (\px,9.6);
\node[red!70!black, font=\Large] at (4.3,9.6) {$\times$};
\node[lbl, above, text=red!70!black] at (\mx-2.2,9.6)
  {health check: connection refused};

\draw[msg] (\px,10.8) -- (\mx,10.8);
\node[lbl, above] at (\px+4.8,10.8)
  {node\_active(reported=primary, pgIsRunning=\textbf{false})};
\draw[rsp] (\mx,11.4) -- (\px,11.4);
\node[lbl, above] at (\mx-2.2,11.4) {goal=primary (unchanged, still deciding)};

\node[note, fill=red!12, draw=red!50, text width=15cm, font=\small]
  at (9.5, 12.9)
  {the keeper's own report already says Postgres is down -- that alone is
   enough to mark the primary unhealthy, whether or not the direct health
   check is also currently failing. (The other path to unhealthy -- the
   keeper goes fully silent \emph{and} the last direct health check was
   bad -- is the Network Partitions scenario, see \S~Monitor can't
   connect to Primary)};

\draw[msg] (\sx,16.0) -- (\mx,16.0);
\node[lbl, above] at (\sx+3.8,16.0) {node\_active(reported=secondary, pgIsRunning=true)};
\draw[rsp] (\mx,16.6) -- (\sx,16.6);
\node[lbl, above] at (\mx-2.2,16.6) {goal=secondary};

\node[note, fill=gray!12, draw=gray!55, text width=13cm, font=\small]
  at (9.5, 17.8)
  {only proceeds because the secondary is itself healthy -- otherwise the
   monitor has no safe candidate and stays put};

%% ── failover begins ──────────────────────────────────────────────────────
\node[seclbl] at (-3.4,19.2) {failover begins, on the primary's own next report};

\draw[msg] (\px,20.1) -- (\mx,20.1);
\node[lbl, above] at (\px+4.8,20.1)
  {node\_active(reported=primary, pgIsRunning=false)};
\draw[rsp, color=pbox, line width=0.8pt] (\mx,20.7) -- (\px,20.7);
\node[lbl, above, text=pbox] at (\mx-2.2,20.7) {goal=draining};

\draw[msg] (\sx,21.9) -- (\mx,21.9);
\node[lbl, above] at (\sx+3.8,21.9) {node\_active(reported=secondary)};
\draw[rsp, color=pbox, line width=0.8pt] (\mx,22.5) -- (\sx,22.5);
\node[lbl, above, text=pbox] at (\mx-2.6,22.5) {goal=prepare\_promotion};

\node[note, fill=sbox!25, draw=sbox!70!black, text width=13cm, font=\small]
  at (9.5, 24.1)
  {secondary catches up on any last WAL still available from the primary,
   then reports success -- or, after
   \texttt{PREPARE\_PROMOTION\_CATCHUP\_TIMEOUT}, reports success anyway and
   the failover continues regardless};

\draw[msg] (\sx,26.2) -- (\mx,26.2);
\node[lbl, above] at (\sx+3.5,26.2) {node\_active(reported=prepare\_promotion)};

%% ── STONITH ──────────────────────────────────────────────────────────────
\node[seclbl] at (-3.4,27.6) {monitor stops the primary (STONITH)};

\draw[msg] (\px,28.5) -- (\mx,28.5);
\node[lbl, above] at (\px+4.6,28.5) {node\_active(reported=draining)};
\draw[rsp, color=pbox, line width=0.8pt] (\mx,29.1) -- (\px,29.1);
\node[lbl, above, text=pbox] at (\mx-2.2,29.1) {goal=demoted};

\draw[lldead] (\px,29.4) -- (\px,\ybot);

\node[note, fill=gray!12, draw=gray!55, text width=13cm, font=\small]
  at (9.5, 30.5)
  {keeper stops (and prevents restarting) Postgres on the old primary --
   guards against a split-brain once the secondary is about to accept
   writes};

%% ── promotion complete ───────────────────────────────────────────────────
\draw[msg, color=pbox, line width=1pt] (\mx,31.7) -- (\sx,31.7);
\node[lbl, above, text=pbox, font=\small\bfseries] at (\mx-2.4,31.7)
  {goal=wait\_primary};

\node[note, fill=green!15, draw=green!50!black,
      font=\small\bfseries, minimum width=8cm]
  at (14, 33.0) {$\checkmark$~~Secondary is the new primary};

%% ── actor footers ───────────────────────────────────────────────────────
\node[actor, fill=gray!30, draw=gray!60, text=gray!30!black]
  at (\px, \ybot+0.35) {\small Primary (demoted)};
\node[actor, fill=pbox, text=ptxt] at (\sx, \ybot+0.35) {New Primary};
\node[actor, fill=mbox, text=mtxt] at (\mx, \ybot+0.35) {Monitor};

\end{tikzpicture}
\end{document}
